\section*{الفصل \ref{ch:g7:heart-and-circulation} --- القلب والدوران}

\begin{solution}{exo:g7:heart-and-circulation:1}
\omterm{def:g7:heart-and-circulation:rooms}{الأذين} الأيمن: يستقبل دم الجسم العائد. \omterm{def:g7:heart-and-circulation:rooms}{والبطين} الأيمن: يضخّه إلى
\omterm{def:g7:how-animals-breathe:four}{الرئتين}. \omterm{def:g7:heart-and-circulation:rooms}{والأذين} الأيسر: يستقبل دم \omterm{def:g7:how-animals-breathe:four}{الرئتين} الغني \omterm{def:g7:what-is-respiration:respiration}{بالأكسجين}.
\omterm{def:g7:heart-and-circulation:rooms}{والبطين} الأيسر: يضخّه حول الجسم كله --- وهو أغلظ الأربع جدارًا.
\end{solution}

\begin{solution}{exo:g7:heart-and-circulation:2}
هي أبواب ذات اتجاه واحد تمنع اندفاع \omterm{prop:g4:journey-of-food:absorb}{الدم} إلى الوراء. و``لَبْ
دَبْ'' انغلاقاتها المزدوجة: فصمّاما الدخول ينغلقان بشدة عند دفع
البطينين (لَبْ)، وصمّاما الخروج ينغلقان عند ارتخائهما (دَبْ).
\end{solution}

\begin{solution}{exo:g7:heart-and-circulation:3}
الدارة الرئوية: من \omterm{def:g7:heart-and-circulation:rooms}{البطين} الأيمن إلى شعيرات \omterm{def:g7:how-animals-breathe:four}{الرئتين} (تحميل
\omterm{def:g7:what-is-respiration:respiration}{الأكسجين} وتفريغ \omterm{def:g7:what-is-respiration:respiration}{ثاني أكسيد الكربون})، عائدةً إلى \omterm{def:g7:heart-and-circulation:rooms}{الأذين} الأيسر.
والدارة الجسمية: من \omterm{def:g7:heart-and-circulation:rooms}{البطين} الأيسر إلى شعيرات الجسم كله (تسليمًا
وجمعًا)، عائدةً إلى \omterm{def:g7:heart-and-circulation:rooms}{الأذين} الأيمن.
\end{solution}

\begin{solution}{exo:g7:heart-and-circulation:4}
\omterm{def:g7:heart-and-circulation:vessels}{الشرايين}: جدران غليظة مرنة --- تحمل \omterm{prop:g4:journey-of-food:absorb}{الدم} بعيدًا عن \omterm{def:g4:heart-and-blood:heart}{القلب} تحت ضغط.
والشعيرات: جدران بسمك \omterm{def:g5:first-look-at-cells:cell}{خلية} واحدة --- وهي موضع كل تبادل. \omterm{def:g7:heart-and-circulation:vessels}{والأوردة}:
واسعة، أرقّ جدارًا، ولها \omterm{def:g7:heart-and-circulation:rooms}{صمامات} ذات اتجاه واحد --- تعيد \omterm{prop:g4:journey-of-food:absorb}{الدم} إلى
\omterm{def:g4:heart-and-blood:heart}{القلب} تحت ضغط منخفض.
\end{solution}

\begin{solution}{exo:g7:heart-and-circulation:5}
لأنه يدفع في الدارة الجسمية --- وهي الحلقة الأكبر بكثير، إلى
أصابع القدم \omterm{def:g5:brain-in-command:system}{والدماغ} وعودةً --- فعليه أن يدفع أشدّ دفع، وجداره
مبني ليطابق ذلك.
\end{solution}

\begin{solution}{exo:g7:heart-and-circulation:6}
في الشعيرات. فشعيرات \omterm{def:g7:how-animals-breathe:four}{الرئتين}: \omterm{def:g7:what-is-respiration:respiration}{أكسجين} إلى الداخل، وثاني أكسيد
كربون إلى الخارج؛ وشعيرات \omterm{prop:g7:digestion-and-nutrients:absorption}{الزغابات}: \omterm{def:g7:digestion-and-nutrients:families}{مغذيات} إلى الداخل؛ وشعيرات
\omterm{def:g3:bones-and-muscles:muscle}{العضلات}: وقود \omterm{def:g7:what-is-respiration:respiration}{وأكسجين} ينزلان، وثاني أكسيد كربون يصعد (وشعيرات
\omterm{def:g1:the-five-senses:sense}{الجلد}: حرارة إلى الخارج).
\end{solution}

\begin{solution}{exo:g7:heart-and-circulation:7}
\omterm{def:g7:heart-and-circulation:rooms}{البطين} الأيمن --- شعيرات \omterm{def:g7:how-animals-breathe:four}{الرئتين} (\omterm{def:g7:what-is-respiration:respiration}{الأكسجين} على متنها) --- \omterm{def:g7:heart-and-circulation:rooms}{الأذين}
الأيسر --- \omterm{def:g7:heart-and-circulation:rooms}{البطين} الأيسر --- \omterm{def:g7:heart-and-circulation:vessels}{شرايين} العنق --- شعيرات \omterm{def:g5:brain-in-command:system}{الدماغ} ---
\omterm{def:g7:heart-and-circulation:vessels}{أوردة} العنق --- \omterm{def:g7:heart-and-circulation:rooms}{الأذين} الأيمن. وتبادل \omterm{def:g5:brain-in-command:system}{الدماغ} يأخذ حصة كبيرة ثابتة
من \omterm{def:g7:what-is-respiration:respiration}{الأكسجين} \omterm{prop:g7:organs-at-work:bill}{والغلوكوز} ويترك \omterm{def:g7:what-is-respiration:respiration}{ثاني أكسيد الكربون} --- وهو أمر القائد
الطلوب الدائم.
\end{solution}

\begin{solution}{exo:g7:heart-and-circulation:8}
\omterm{prop:g4:heart-and-blood:pulse}{النبض} هو تمدّد جدار \omterm{def:g7:heart-and-circulation:vessels}{الشريان} المرن عند كل دفقة بطينية --- ضغط عالٍ
يلقى جدارًا غليظًا مطاطيًّا. أما عند \omterm{def:g7:heart-and-circulation:vessels}{الأوردة} فقد أُنفقت الدفقة في
عبور الشعيرات: ضغط منخفض ناعم في أنبوب واسع ليّن --- فلا شيء بقي
لينبض.
\end{solution}

\begin{solution}{exo:g7:heart-and-circulation:9}
العودة الصاعدة: فدم \omterm{def:g7:heart-and-circulation:vessels}{الأوردة} من \omterm{def:g1:my-body:parts}{الرجلين} عليه أن يصعد ضد الجاذبية
تحت ضغط منخفض، \omterm{def:g7:heart-and-circulation:rooms}{والصمامات} تحفظ كل مكسب. \omterm{def:g3:bones-and-muscles:muscle}{وعضلات} \omterm{def:g1:plants-around-us:parts}{ساق} الجندي الواقف
بلا حراك لا تعصر \omterm{def:g7:heart-and-circulation:vessels}{الأوردة} قط، فيتعثّر الصعود، ويتجمع \omterm{prop:g4:journey-of-food:absorb}{الدم} في
\omterm{def:g1:my-body:parts}{الرجلين} وينخفض تموين \omterm{def:g5:brain-in-command:system}{الدماغ} --- ومن هنا الإغماء. فالأقدام السائرة
مضخّات \omterm{def:g7:heart-and-circulation:vessels}{أوردة}.
\end{solution}

\begin{solution}{exo:g7:heart-and-circulation:10}
\omterm{def:g7:heart-and-circulation:vessels}{شريان} داخل من الدارة الجسمية؛ وفي الداخل تلتقط شعيرات \omterm{prop:g7:digestion-and-nutrients:absorption}{الزغابات}
\omterm{def:g7:digestion-and-nutrients:families}{مغذيات} الوجبة
(\cref{prop:g7:digestion-and-nutrients:absorption})؛ والوصلة
الخاصة: أن وريده يمضي أولًا إلى \omterm{prop:g7:digestion-and-nutrients:absorption}{الكبد} للفرز والإيداع، وعندئذ فقط
ينضمّ \omterm{prop:g4:journey-of-food:absorb}{الدم} إلى العودة العامة إلى \omterm{def:g7:heart-and-circulation:rooms}{الأذين} الأيمن. فالرسم يعيد حكاية
فصل \omterm{def:g4:journey-of-food:digestion}{الهضم}.
\end{solution}

\begin{solution}{exo:g7:heart-and-circulation:11}
دم الغرف يندفع عابرًا؛ أما جدران \omterm{def:g4:heart-and-blood:heart}{القلب} الغليظة نفسها، وهي تعمل بلا
راحة، فتحتاج إلى تسليم شعيري خاص بها كأي \omterm{def:g3:bones-and-muscles:muscle}{عضلة} --- فتتفرع \omterm{def:g7:heart-and-circulation:vessels}{شرايين}
تغذية صغيرة على \omterm{def:g4:heart-and-blood:heart}{القلب} نفسه. وهي إذا كُسيت خنقت خط وقود المضخة
نفسها: فالعضو الوحيد الذي لا يستطيع أن يتوقف هو الذي يجوّعه
انسدادها.
\end{solution}

\begin{solution}{exo:g7:heart-and-circulation:12}
كل تبادل يعيش الجسم به --- \omterm{def:g7:what-is-respiration:respiration}{الأكسجين} عند \omterm{def:g7:lungs-and-gas-exchange:alveoli}{الأسناخ}، \omterm{def:g7:digestion-and-nutrients:families}{والمغذيات} عند
\omterm{prop:g7:digestion-and-nutrients:absorption}{الزغابات}، والوقود والفضلات عند كل \omterm{def:g3:bones-and-muscles:muscle}{عضلة} --- لا يحدث إلا عبر جدران
الشعيرات؛ \omterm{def:g7:heart-and-circulation:vessels}{والشرايين} \omterm{def:g7:heart-and-circulation:vessels}{والأوردة} لا تفعل إلا أن تجلب \omterm{prop:g4:journey-of-food:absorb}{الدم} إلى تلك
الجدران وتأخذه بعيدًا، \omterm{def:g4:heart-and-blood:heart}{والقلب} لا يفعل إلا أن يُبقيه متحركًا.
فالمضخة والطرق السريعة موجودة من أجل الممرات الرفيعة كالشعر: أوقف
حركة الشعيرات في أي مكان يمت ذلك العضو، مهما تكن السباكة سليمة.
\end{solution}

\begin{solution}{pb:g7:heart-and-circulation:1}
\textbf{1.} \omterm{prop:g4:heart-and-blood:pulse}{النبض}: عدّ دفقات \omterm{def:g7:heart-and-circulation:vessels}{الشرايين} --- وتيرة المضخة مقروءةً عند
\omterm{def:g1:my-body:joint}{رسغ}. وضغط \omterm{prop:g4:journey-of-food:absorb}{الدم}: ضغط \omterm{def:g7:heart-and-circulation:vessels}{الشريان} عند الدفقة وبينها --- حمل الطرق
السريعة. و``لَبْ دَبْ'' في السمّاعة: \omterm{def:g7:heart-and-circulation:rooms}{الصمامات} مسموعةً وهي تنغلق ---
سلامة الأبواب.

\textbf{2.} \omterm{prop:g4:heart-and-blood:pulse}{نبض} الراحة المنخفض عند مريضة لائقة يعني بطينًا قويًّا
يحرّك دمًا أكثر في الدقّة: \omterm{def:g4:heart-and-blood:heart}{فالقلب} المدرَّب يتباطأ هادئًا، والاحتياط
في يده.

\textbf{3.} أبواب \omterm{def:g7:heart-and-circulation:rooms}{الصمامات}: صمّاما الدخول عند دفع البطينين (لَبْ)،
وصمّاما الخروج عند ارتخائهما (دَبْ) --- أربعة أبواب تنغلق في
زوجين، انغلاقًا نظيفًا.

\textbf{4.} \omterm{def:g7:heart-and-circulation:rooms}{الأذين} الأيمن، \omterm{def:g7:heart-and-circulation:rooms}{فالبطين} الأيمن، فشعيرات \omterm{def:g7:how-animals-breathe:four}{الرئتين}،
\omterm{def:g7:heart-and-circulation:rooms}{فالأذين} الأيسر، \omterm{def:g7:heart-and-circulation:rooms}{فالبطين} الأيسر، \omterm{def:g7:heart-and-circulation:vessels}{فشرايين} الجسم، فشعيرات الجسم
(\omterm{def:g3:bones-and-muscles:muscle}{عضلة} مثلًا)، \omterm{def:g7:heart-and-circulation:vessels}{فالأوردة} --- وعودةً إلى \omterm{def:g7:heart-and-circulation:rooms}{الأذين} الأيمن.

\textbf{5.} كل دقّة تفقد جزءًا من دفقتها راجعًا عبر الرشح، فتصغر
توصيلات الدارة الجسمية في الدقّة؛ ويعوّض \omterm{def:g4:heart-and-blood:heart}{القلب} بدقات أكثر وعمل
أشقّ، فتقلّ قدرة الاحتمال --- واللغط هو التوصيل ضائعًا ارتدادًا.

\textbf{6.} ضغط \omterm{prop:g4:journey-of-food:absorb}{الدم} مرتفع (فالطرق السريعة المتصلّبة الضيقة تقاوم
الدفقة) وتمدّد \omterm{def:g7:heart-and-circulation:vessels}{الشريان} في \omterm{prop:g4:heart-and-blood:pulse}{النبض} متغيّر؛ \omterm{def:g7:heart-and-circulation:rooms}{والبطين} الأيسر، وهو يدفع
كل جولة في مقابل دارة متصلّبة، يعمل ساعات إضافية ويغلظ.

\textbf{7.} الكسوة في \omterm{def:g7:heart-and-circulation:vessels}{شرايين} \omterm{def:g4:heart-and-blood:heart}{القلب} المغذية تخنق خط وقود \omterm{def:g3:bones-and-muscles:muscle}{العضلة}
التي لا تستطيع أن تستريح أبدًا --- وهو الطريق الكلاسيكي إلى
احتشاء \omterm{def:g4:heart-and-blood:heart}{القلب}.

\textbf{8.} \omterm{def:g7:heart-and-circulation:rooms}{صمامات} \omterm{def:g7:heart-and-circulation:vessels}{الأوردة} ذات الاتجاه الواحد ومضخة العضل الغائبة
مع انعدام الحركة: فمع الوقوف بلا حراك لا تعصر \omterm{def:g1:plants-around-us:parts}{السيقان} قط، فتتعثّر
مكاسب \omterm{def:g7:heart-and-circulation:rooms}{الصمامات}، ويتجمع \omterm{prop:g4:journey-of-food:absorb}{الدم} في الكاحلين. والمشي يعصر \omterm{def:g7:heart-and-circulation:vessels}{الأوردة}
إيقاعيًّا --- فالراحة ميكانيكية.

\textbf{9.} لأن جولة الجسم تنفق \omterm{def:g7:what-is-respiration:respiration}{أكسجين} \omterm{prop:g4:journey-of-food:absorb}{الدم}: ولا تستطيع إعادة
تحميله إلا الدارة الرئوية. وتخطّي \omterm{def:g7:how-animals-breathe:four}{الرئتين} سيرسل دمًا منفَقًا في
توصيلات بشاحنة فارغة --- فالحلقة المزدوجة تزوّد ثم تسلّم، ويفرض
ذلك جانبا \omterm{def:g4:heart-and-blood:heart}{القلب}.

\textbf{10.} \omterm{def:g7:heart-and-circulation:vessels}{الشريان}: غليظ مرن عالي الضغط --- طريق سريع خارج.
والشعيرة: بسمك \omterm{def:g5:first-look-at-cells:cell}{خلية} واحدة، تتخلل كل شيء --- وهي ممرات التسليم حيث
تحدث كل التبادلات. \omterm{def:g7:heart-and-circulation:vessels}{والوريد}: واسع ذو \omterm{def:g7:heart-and-circulation:rooms}{صمامات} منخفض الضغط --- وهو
العودة.

\textbf{11.} كل عطل في الاستشارات همّ بالضبط حيث مسّ خدمة
الشعيرات: فالرشح جوّع التوصيلات، والكسوة خنقت خطوط التموين،
والتجمع عطّل العودات. فالقراءات \omterm{def:g7:heart-and-circulation:rooms}{والصمامات} والطرق السريعة وسائل؛
وحركة الشعيرات --- أي التبادلات --- هي الغاية.

\textbf{12.} تحرّك كل يوم --- فتقوى المضخة وتبقى الطرق السريعة
لدنة؛ وكُلْ متوازنًا --- فتبقى جدران \omterm{def:g7:heart-and-circulation:vessels}{الشرايين} نقية من الكسوة؛ ولا
دخان --- فلا تضيق \omterm{def:g4:heart-and-blood:vessels}{الأوعية} ولا تتصلّب، وتبقى مقاعد \omterm{prop:g4:journey-of-food:absorb}{الدم} شاغرة
\omterm{def:g7:what-is-respiration:respiration}{للأكسجين}.
\end{solution}
